\documentclass[sigconf,nonacm]{acmart}

% codoc — UIST short paper (2 pages).
% Draft. Citations are written inline in APA style; the reference list at the
% end is APA-formatted. BibTeX to be added later.

\settopmatter{printccs=false, printacmref=false}
\setcopyright{none}
\pagestyle{plain}

\begin{document}

\title{codoc: Reading and Changing a Codebase Through an Editable Map of Its Intent}

\author{Anonymous}
\affiliation{%
  \institution{Anonymous Institution}
  \country{}
}

\begin{abstract}
As AI agents write and edit a growing share of code, the account of what the code
is for --- why a piece of code exists, and how one feature is spread across many
files --- falls behind, because the agent maintains files but not that account,
and the documents meant to hold it are authored to feed the agent context rather
than kept faithful to the code. codoc treats the documentation of a codebase as a
shared representation that a person and an agent both read and edit. It is built
from the code: the repository is parsed into code chunks, each chunk is labeled,
and the labels are clustered into a hierarchy of features, so every section of the
document corresponds to a definite slice of the code and that correspondence is
kept live. A person reads the document to understand the system and edits it to
say what should change; the agent implements the change and surfaces the result
back into the document as a diff to review; and both sides update incrementally,
so no one has to re-read a regenerated document. We describe the design through a
running example --- making a small coding agent's model calls resilient to a flaky
server --- and its implementation as a Visual Studio Code extension.
\end{abstract}

\keywords{program comprehension, code navigation, AI coding agents, intent, literate programming}

\maketitle

\section{Introduction}

Managing a codebase over its lifetime demands a clear picture of not just what it
does but how it works. To plan a change, review one, or decide where new code
belongs, a developer needs to know how the parts fit together, not only what the
system delivers on the surface. Coding through LLM agents leaves that picture
thin: the agent conveys what the code does at a high level, but the lower-level
account of how it works --- the account that informs each next decision in
iteration and maintenance --- is not something it keeps.

The usual remedy is a companion document kept alongside the code: a
\texttt{CLAUDE.md}, a set of Cursor rules, a memory bank. These soon turn into
long change logs, useful as context for the agent but of little help to a person
trying to understand the system. The root problem is authorship: no human wrote
them. They are the agent's summaries of code the person neither wrote nor
condensed, so a developer can neither trust them as an account of the system nor
edit them to specify what should change next.

codoc instead builds the document bottom-up from the code, so it stays faithful to
the codebase and serves as a \emph{shared representation} that a person and an
agent co-maintain. The repository is parsed into an
intermediate representation of code chunks; each chunk is assigned a semantic
label; and the labels are clustered into a hierarchy of \emph{features}, named
units of intent. Each section of the document therefore corresponds to a definite
slice of the code, and the correspondence is kept live rather than computed once:
a change to any code in a feature's slice, made by a person or by an agent, is
detected and surfaced in that section.

On top of this mapping codoc runs an editing loop that resembles collaborative
document editing. A person edits the document to say what should change; the agent
picks up the implied task, implements it in code, and surfaces the resulting
changes --- including the ones the person did not spell out --- back into the
document as a diff to review. Because both the document and the code are updated
incrementally, the person never confronts a regenerated document they must
re-read from scratch; they see only what changed and who changed it. Our
contribution is this design: one editable representation of a codebase's intent,
kept faithful to the code in both directions, through which a person comprehends
the system, directs the agent, and reviews its work. We describe it through a
running example and its implementation as a Visual Studio Code extension.

\section{Representing intent: a recurring goal, and why it kept failing}

Describing a program in terms of its intent is not a new wish. Knuth (1984) argued
that a program should be written to explain to a person what we want the computer
to do, with the machine as the second reader. His approach kept the prose and the
code consistent at the moment they were written, from a single source. It said
nothing about what happens when the code later changes under maintenance, and in
practice the explanation went stale or the code became hard to change once it was
embedded in prose. The method also assumed one author keeping both sides aligned
by hand.

Understanding code has long been framed as mapping human concepts onto the code
that implements them. Biggerstaff et al.\ (1993) called this the concept
assignment problem and described code as carrying both an informal expression of
intent for people and a formal one for the machine. Rajlich and Wilde (2002) put
it in terms that match daily work: a change is stated in domain concepts, the
concept is known but where it lives in the code is not, and one concept is usually
spread across many files. People built tools to find that code, surveyed by Dit et
al.\ (2013), but these ran once per task and returned a ranked list that a
developer read and discarded. The mapping was recomputed each time, never kept.

Others tried the opposite: make intent the source and generate code from it.
Simonyi et al.\ (2006) stored intent as the primary artifact and generated code
from it, but the surface was a specialized structured editor. Voelter et al.\
(2014) documented the cost of that style of editor: people had to edit the syntax
tree rather than type freely, ordinary tools like diff and version control no
longer worked, and the effort of getting used to it was high.

The same goal returns in today's agent tooling, in three forms. One keeps a
natural-language artifact beside the code to steer the agent: GitHub Next's
SpecLang (2023) writes a program as a prose specification the model compiles to
code, while Cursor rules (2024) and Cline's memory bank (2025) persist project
conventions and state for the agent to reload on each turn. These make the
agent's output more controllable, but the artifact is authored for the model and
is not bound to the code it governs, so it drifts exactly as earlier
documentation did. A second form conditions code generation on documentation
directly: DocPrompting (Zhou et al., 2023) retrieves library docs to generate
calls to functions the model has not seen, and DocCGen (Pimparkhede et al., 2024)
constrains generation against structured documentation to raise syntactic
correctness and controllability. A third builds an explicit intermediate
structure before generating: RPG (Microsoft Research, 2025) plans a repository as
a graph of capabilities and data flows and generates the code from it. All three
improve how code is produced \emph{from} a description, but the description is an
input consumed once, not a living representation a person keeps reading and
editing as the code evolves.

Across these efforts the same goal recurs --- to hold a codebase in terms of its
intent --- and two obstacles defeat it. Keeping the map between intent and code
current was manual, so it rotted; and the editors that made intent the primary
artifact were too heavy to adopt. Both obstacles have since moved. An agent can
now read code and re-derive what each part does, and write code from a stated
intent, cheaply enough to run continuously. codoc builds on this: it derives the
intent-to-code map automatically and keeps it current in both directions, and it
keeps the authoring surface as plain text so a person never leaves the editor
they already use.

\section{A running example, and what codoc shows}

Maya has taken over a small coding agent: a few thousand lines that drive a
read-think-act tool-use loop against a locally hosted model. The original author
has left, and a user reports that the local model server is flaky --- a single
dropped request aborts the whole session. Maya has to make model calls resilient,
in code she has not read.

She opens \texttt{tree.codoc} in the Codoc Tree editor, the view codoc keeps for
this repository (Figure~\ref{fig:tree}). The outline is a projection of codoc's
store, and it gives her the shape of the system in about a minute: the agent is a
loop orchestrator that builds a size-capped prompt, parses the model's tool calls,
dispatches them through a validated tool registry, and runs file and shell tools
behind a path sandbox, while the model itself is reached through a separate
\emph{model backend clients} area. Each feature links to the exact code that
implements it, so she clicks the link on ``Ollama model backend client'' and lands
on the \texttt{OllamaModelClient} class that makes the HTTP call. She has run no
search and opened no unrelated file.

\begin{figure}[t]
\footnotesize
\begin{verbatim}
- Agent runtime core
  - Agentic tool-use loop orchestrator
      Drives the read-think-act loop: sends prompts,
      parses replies, executes tools, spawns sub-agents.
      [MiniAgent](codoc:mini_coding_agent.py#MiniAgent)

- Model backend clients
  - Ollama model backend client
      > On a transient HTTP error or timeout, retry with
      > **exponential backoff** (max 3 attempts) instead
      > of aborting the run. Add a test for the retry path.
      Talks to a local Ollama /api/generate endpoint for
      completions.
      Runbook: [flaky local inference](wiki/runbook/ollama)
      [OllamaModelClient](codoc:mini_coding_agent.py#OllamaModelClient)
\end{verbatim}
\caption{A fragment of \texttt{tree.codoc} after Maya's edits. Features nest under
an area and each links to the code that implements it (feature ids are hidden in
the editor). The two quoted lines are a steering note to the agent, the bold span
marks the requirement, and the plain link points the agent at a document to read
first.}
\label{fig:tree}
\end{figure}

\section{Editing the intent instead of changing the code directly}

Reading the client, Maya confirms the cause: the call to \texttt{/api/generate}
sets a timeout but never retries, so one transient failure propagates up and ends
the run. She states the fix where the feature is described, not in a chat window.
On ``Ollama model backend client'' she adds two quoted note lines --- on a
transient HTTP error or timeout, retry with exponential backoff up to three times
instead of aborting; and add a test for the retry path --- marks ``exponential
backoff'' in bold so the agent treats it as the requirement, and pastes a link to
the runbook for the flaky server so the agent reads it first. codoc does not treat
these as plain prose: quoted lines become steering directives, the bold span is
the focus of the change, and the external link is fetched before implementation.
The instruction stays attached to the feature it concerns, with a place to live
and a version, rather than scattered across throwaway prompts (Prompts Are
Programs Too!, 2024).

The edits are held as a draft until Maya hands them to the agent, which starts
Loop~B: the tree-to-code direction. The agent reads the runbook, wraps the
completion call in retry-with-backoff, and adds the test. The change returns as a
normal diff, with the agent's edits shown as tracked changes she accepts or
rejects inline, together with the updated links from the feature to the new code.
She reviews and accepts.

The retry logic is a new helper function that no existing feature owns. Because
the code no longer matches the outline, codoc's other direction --- Loop~A, which
watches the code --- surfaces this as a proposal rather than an edit made behind
her back: a new child feature under the model client for the retry helper,
rendered in place as a ghost row with inline Accept and Reject. Maya accepts it,
because it matches how she now understands the code. When she is done the outline
is correct, the fix is in, and the reason for it --- down to the runbook link and
the bolded requirement --- is recorded where the code is described, not in a
ticket that will be closed and forgotten.

\section{How codoc keeps the two in sync}

We make the two directions precise. The repository is indexed into a set of
AST-level chunks $\mathcal{C}$; each chunk $c = (\mathit{file}, \mathit{sym}, h)$
carries its location and a content hash $h$ over its normalized tokens. The
outline is a feature tree $\mathcal{T} = (\mathcal{F}, \prec)$ with a forest
parent relation $\prec$, and attribution is a partial map
$\mathcal{B} \colon \mathcal{C} \to \mathcal{F} \cup \{\bot\}$ that binds each
chunk to at most one feature --- a uniqueness constraint on
$(\mathit{file}, \mathit{sym})$ --- while one feature owns many chunks across many
files. This map is the correspondence between a section of the document and its
slice of the code.

\textbf{Loop~A reflects code into the tree.} A pass diffs the fresh index against
$\mathcal{B}$ to compute a changeset
$\Delta = (\Delta^{+}, \Delta^{-}, \Delta^{\sim})$ of chunks added, removed, and
modified (bound, but $h$ no longer matches the stored fingerprint). Every change
it can resolve mechanically it resolves without a model: a modified chunk
refreshes its binding, a removed chunk detaches, and a chunk whose hash reappears
at a new key re-attaches to its prior feature, carrying attribution across moves
and renames. Only the residual --- genuinely new code, or a feature that lost its
last binding --- goes to a single LLM call per pass, which sees $\Delta$, a
locality subtree of $\mathcal{T}$, and \emph{every} feature title. Supplying the
full title set to one call is what prevents duplicate or split features, with no
post-hoc deduplication. Safe operations (binding maintenance, small prose
amendments) apply automatically; structural ones (add, move, retire) become
proposals a person accepts or rejects. The direction is workable because code that
people write is regular enough that a model can reliably tell which code serves
which feature (Hindle et al., 2012); we never need an exact, reversible grammar
between outline and code, the requirement that made earlier round-trip tools
brittle once code was hand-edited (Siegl, 2015).

\textbf{Loop~B realizes tree edits into code.} Each applied edit is classified by
whether it implies a code change. Descriptive prose merely persists; imperative
prose, an accepted plan node, or a steering note mints a \emph{directive} --- a
work order assembled from the feature's description, its bound symbols, and its
\texttt{Consult} links --- and queues it for the live coding session. The session
implements the directive, reflects the new code back, and Loop~A closes the cycle.

\textbf{Conflicts resolve by a deterministic priority, not a merge heuristic.}
Let $H \subseteq \mathcal{F}$ be the \emph{hold set}: features with a pending
document-side intent, whether an unsettled suggestion or a queued directive. While
$f \in H$, code-side intent operations on $f$ are suppressed,
\[
  \mathrm{suppress}(o, H) \iff
  o.\mathit{kind} \in \{\textsf{amend}, \textsf{retire}, \textsf{move}\}
  \;\wedge\; o.f \in H,
\]
while binding maintenance ($\textsf{attach}$, $\textsf{detach}$,
$\textsf{refresh}$) is never suppressed --- bindings are attribution, not intent,
and must stay correct regardless of who is mid-edit. The rule is that the side
expressing \emph{intent} (the document) outranks the side that only
\emph{observes} (the code reflector), so a person's in-flight edit is never
clobbered.

\section{Status and next steps}

codoc runs today on real repositories as a Visual Studio Code extension. Reading
the outline, editing it, directing the agent, and reviewing structural proposals
work end to end. The rough edges are honest ones: the model pass adds latency,
proposals on large refactors are not always right, and the first outline built
from a cold start needs cleanup.

Three directions follow. First, \emph{shared editing.} A single outline that a
team reads and edits is a natural place for more than one person, and for people
and agents, to work at once; we are moving the store's projection onto a CRDT
(pycrdt) so concurrent edits from several authors and the reflecting loop merge
without a lock or a lost write. Second, \emph{more media than prose.} A feature is
already a typed, bindable node, and the same lift-and-lower machinery applies to
non-prose media --- a dependency diagram lifted from the code graph, a screenshot
or a linked document consulted before implementation --- so the outline can carry
the representation that fits each feature rather than prose alone. Third,
\emph{evaluation.} We have not yet run a user study; the question we most want to
answer is whether the outline holds up as the primary place people work over
months of real change, not just for a single fix.

\begin{thebibliography}{99}

\bibitem{allamanis2018}
Allamanis, M., Barr, E. T., Devanbu, P., \& Sutton, C. (2018). A survey of machine
learning for big code and naturalness. \textit{ACM Computing Surveys, 51}(4),
1--37.

\bibitem{biggerstaff1993}
Biggerstaff, T. J., Mitbander, B. G., \& Webster, D. (1993). The concept assignment
problem in program understanding. In \textit{Proceedings of the 15th International
Conference on Software Engineering} (pp. 482--498). IEEE.

\bibitem{cline2025}
Cline. (2025). \textit{Memory Bank} [Software documentation]. Retrieved from
https://docs.cline.bot

\bibitem{cursor2024}
Cursor (Anysphere). (2024). \textit{Rules for AI} [Software documentation].
Retrieved from https://docs.cursor.com/context/rules

\bibitem{dit2013}
Dit, B., Revelle, M., Gethers, M., \& Poshyvanyk, D. (2013). Feature location in
source code: A taxonomy and survey. \textit{Journal of Software: Evolution and
Process, 25}(1), 53--95.

\bibitem{speclang2023}
GitHub Next. (2023). \textit{SpecLang} [Research prototype]. Retrieved from
https://githubnext.com/projects/speclang

\bibitem{hindle2012}
Hindle, A., Barr, E. T., Su, Z., Gabel, M., \& Devanbu, P. (2012). On the
naturalness of software. In \textit{Proceedings of the 34th International
Conference on Software Engineering} (pp. 837--847). IEEE.

\bibitem{knuth1984}
Knuth, D. E. (1984). Literate programming. \textit{The Computer Journal, 27}(2),
97--111.

\bibitem{rpg2025}
Microsoft Research. (2025). RPG: A repository planning graph for unified and
scalable codebase generation. \textit{arXiv:2509.16198}.

\bibitem{doccgen2024}
Pimparkhede, S., Kammakomati, M., Tamilselvam, S., Kumar, P., Kumar, A. P., \&
Bhattacharyya, P. (2024). DocCGen: Document-based controlled code generation. In
\textit{Proceedings of the 2024 Conference on Empirical Methods in Natural
Language Processing (EMNLP)}. ACL.

\bibitem{promptsprograms2024}
[Authors to confirm]. (2024). Prompts are programs too! Understanding how
developers build software containing prompts. \textit{arXiv:2409.12447}.

\bibitem{rajlich2002}
Rajlich, V., \& Wilde, N. (2002). The role of concepts in program comprehension. In
\textit{Proceedings of the 10th International Workshop on Program Comprehension}
(pp. 271--278). IEEE.

\bibitem{siegl2015}
Siegl, D. (2015). Why round-trip engineering does not work. LieberLieber (blog).

\bibitem{simonyi2006}
Simonyi, C., Christerson, M., \& Clifford, S. (2006). Intentional software. In
\textit{Proceedings of OOPSLA 2006} (pp. 451--464). ACM.

\bibitem{voelter2014}
Voelter, M., Siegmund, J., Berger, T., \& Kolb, B. (2014). Towards user-friendly
projectional editors. In \textit{Software Language Engineering (SLE 2014)}, LNCS
8706 (pp. 41--61). Springer.

\bibitem{docprompting2023}
Zhou, S., Alon, U., Xu, F. F., Jiang, Z., \& Neubig, G. (2023). DocPrompting:
Generating code by retrieving the docs. In \textit{The Eleventh International
Conference on Learning Representations (ICLR)}.

\end{thebibliography}

\end{document}
